% 1.4 — Objectifs
\section{Objectifs du projet}
\label{sec:objectifs}

Les objectifs découlent de la problématique ; chacun est formulé de façon à pouvoir être
vérifié dans le code et les tests.

\begin{table}[htbp]
\centering
\caption{Objectifs du stage et critère de vérification de chacun}
\label{tab:objectifs}
\small
\begin{tabularx}{\linewidth}{@{}L{0.7cm}L{5.6cm}Y@{}}
\toprule
& \textbf{Objectif} & \textbf{Comment on le vérifie} \\
\midrule
O1 & \textbf{Prédire} l'effet d'une ligne à partir de sa structure, pas de son texte. &
  Chaque exemple de la table de règles (\nbexemples{}) est un test qui compare l'impact
  prédit à l'impact attendu. \\
O2 & \textbf{Ne jamais autoriser ce qui n'est pas lu.} Une ligne illisible n'obtient pas
  ALLOW. & Tests M2.3a : une ligne « non comprise » va au filet, à un point de contrôle,
  sinon elle est escaladée. \\
O3 & \textbf{Ne jamais revendiquer une récupération non prise.} &
  Tests des captures partielles (tout ou rien), de V17 (point de contrôle hors projet). \\
O4 & \textbf{Voir pendant l'exécution} ce que la lecture ne peut pas prédire. &
  Garde d'appels système (Linux) et montage WinFsp (Windows) testés sur les deux systèmes. \\
O5 & \textbf{Traiter le distant honnêtement} : distinguer « demander à la personne » et
  « refuser ». & Table \fichier{policy.ESCALATION\_CHANNEL} ; une surface sans ligne est
  refusée. \\
O6 & \textbf{Les deux systèmes} : Linux et Windows, avec les mêmes garanties. &
  Suite complète passée sur les deux systèmes avant chaque atterrissage. \\
O7 & \textbf{Ne rien rendre plus permissif sans le vouloir.} &
  Portail de décision : \nbcorpus{} commandes comparées avant et après chaque partie. \\
O8 & \textbf{Résister à l'agent} : il ne doit pas pouvoir désarmer l'outil. &
  Plan de contrôle refusé à toutes les portes ; tâche Windows durcie (W-SEC). \\
\bottomrule
\end{tabularx}
\end{table}

Un objectif est volontairement \emph{absent} : faire de \demo{} une barrière contre un
adversaire déterminé. L'outil protège contre les \emph{erreurs} d'un agent coopératif ;
ce modèle de menace, hérité de la bêta et affiché par \cmd{demo\_cli doctor}, est discuté
à la section~\ref{sec:menace}.
